Draining Tank Time

QUASI-STEADY DRAINAGE
Tank level changes slowly, but outlet is modeled by instantaneous steady energy.

IDEAL LARGE TANK with outlet Ao:
$$V_o\approx\sqrt{2gh}$$
$$Q=A_o\sqrt{2gh}$$

TANK CONTINUITY
$$A_t(h)\frac{dh}{dt}=-Q(h)$$

CONSTANT AREA At, no losses:
$$t=\frac{2A_t}{A_o\sqrt{2g}}$$
$$\times(\sqrt{h_i}-\sqrt{h_f})$$

If discharge coefficient Cd specified, replace Ao by Cd*Ao.
If pipe/fitting loss K specified instead, derive Q(h) from full energy; DO NOT automatically also add Cd.

IF AREA VARIES WITH HEIGHT
Separate variables with At(h) inside the integral:
$$t=\int_{h_f}^{h_i}\frac{A_t(h)}{Q(h)}dh$$

DISTINGUISH PROBLEM TYPES
Acceleration / pressure startup -> unsteady Bernoulli with local inertia integral.
Time to drain -> energy for Q(h) + changing tank level continuity.

SIGN CHECK
During drainage dh/dt < 0 but elapsed t > 0. Reverse limits properly.
